\subsection{Complexity analysis\label{sec:comparison to state of the art}}
In this section we give a detailed complexity analysis of our collisionless quantum Boltzmann method (CQBM) algorithm. 
Since the most competitive quantum hardware can only implement single- and two-qubit gates natively \cite{Rigetti} \cite{IBM}, we decompose the multi-qubit operations into two-qubit gates to find a realistic quantum complexity for our method. Specifically we choose to give the complexity in terms of the CNOT gates, since these are native to IBM QPU. Though other vendors' QPUs have different native gate sets, e.g., Rigetti \cite{Rigetti}, the restriction to only one- and two-qubit gates holds for nearly all technologies of today. 

We subsequently compare our found complexity to the complexity for the current best known quantum algorithm for the collisionless Boltzmann equation \cite{Todorova2020}. In doing so we show that our method not only outperforms the current  by implementing the quantum reflection step in a fail-safe manner, but we also reach a better complexity.

\subsection{Complexity of multi-controlled NOT operations\label{ssec:multinot compl}}
In order to provide a cost analysis for our algorithm, we first need to determine the cost associated to implementing a multi-controlled NOT gate in terms of natively implementable 2-qubit gates. Specifically we will give the complexity of implementing multi-controlled NOT operations when decomposed in CNOT gates. While CNOT is one of the native gates for IBM QPUs \cite{IBM} it needs to be emulated by other two-qubit gates on other quantum hardware such as the QPUs from Rigetti \cite{Rigetti}. However, in these cases the CNOT gate can be decomposed using a fixed number of the natively implemented two-qubit gates, so this does not effect the overall complexity. 
Multiple ways of decomposing multi-controlled NOT gates into CNOT gates are known in literature. We distinguish several cases, one without ancilla qubits, one with a linear amount of ancilla qubits and multiple with a constant amount of ancilla qubits. \\

The method to decompose multi-controlled NOT gates into two-qubit operations without any ancilla qubits was given by Barenco et al. in \cite{Barenco1995}. The reported decomposition requires $48(p+1)^2 + \Theta(p)$ `basic operations'\footnote{Here the term basic operations implies either a CNOT or single qubit gate.}. From this we can deduce that approximately $c p^2$ with $22<c<48$ CNOT gates will be required for this method with complexity of $\mathcal{O}\left (p^2\right )$ CNOT gates. \\

With a linear amount of ancilla qubits we can use a more efficient decomposition in terms of CNOT gates required as described by Nielsen and Chuang in \cite{mikeandike}. For a multi-controlled operation with $p$ control qubits $p-1$ ancillae and $2(p-1)$ Toffoli gates\footnote{A Toffoli gate is a CCNOT gate, i.e., a controlled NOT gate with two control qubits.} are required. Each Toffoli gate in turn will be decomposed using 6 CNOT gates (and some other operations) \cite{mikeandike}. Therefore a total of $12(p-1)$ CNOT gates is required in combination with $p-1$ ancilla qubits to decompose a C$^p$NOT operation into native one- and two-qubit gates. \\

We can, however, also choose to use a constant amount of ancillae. 
The paper by Barenco et al. \cite{Barenco1995} provides a method to decompose C$^p$NOT operations for $p\geq 5$ using a single ancilla, $8(p-3)$ Toffoli gates and $ 48 (p+2)$ `basic operations' leading to $c p$ CNOT gates required to implement this method with $70 \leq c \leq 96$.

Another method we could use to implement a multi-control NOT gate using a single ancilla is by making use of the comparison operation. One can then find whether or not the $p$ control qubits are in the state $\ket{q_1 \dots q_k} = \ket{1 \dots 1}$ as required. This is done by simply using the comparison operation to check whether $q \geq 2^p-1$ holds, where $q$ is the integer value encoded by the qubits $q_1 \dots q_p$. The cost of such a comparison operation depends on the costs of the constant addition operation. Using the Draper adder as described in Section \ref{sec:efficient quantum convection operation} this method requires $4p^2 + 6p + 3$ CNOT operations\footnote{Since a QFT operation applied to $k$ qubits consists of the following two-qubit gates, first $\frac{k(k-1)}{2}$ controlled phase shift gates followed by $\lfloor k/2 \rfloor$ swap gates. A swap gate can be decomposed using 3 CNOT operations and a controlled phase shift operation requires two CNOT operations. Therefore the total cost of a QFT operation applied to $k$ qubits is $k(k-1) + \lfloor \frac{3k}{2} \rfloor$ CNOT operations. Then the total costs given in the text follows from the fact that the Draper Adder implemented comparison operation consists of first a QFT and QFT$^\dagger$ applied to $p+1$ qubits followed by  a QFT and QFT$^\dagger$ applied to $p$ qubits.}. Therefore this method requires $\mathcal{O} \left ( p^2\right )$ CNOT gates and one additional ancilla.

This implementation appears to match the complexity of the `recursion method', provided by Qiskit \cite{ Qiskit}, which also requires 1 ancilla qubit. 
IBM does not, however, give a complexity analysis in terms of the amount of CNOT gates required or a reference detailing the implementation. We numerically analyzed the complexity of this implementation by running it for several values of $p$ (up to $p=40)$ and found experimentally that their implementation requires $\approx 2 p^2$ CNOT gates. \\

Having seemingly the best trade-off between the amount of CNOT gates and the extra ancilla qubits required for realistic values of $p$, we choose to use the `comparison' or `recursion' multi-controlled NOT decomposition to analyze the complexity of C$^p$NOT gates in terms of CNOT gates. On top of that our numerical analysis (for values of $p$ up to 40) shows it is the complexity of the decomposition method implemented in the  quantum SDK Qiskit \cite{Qiskit}.
And so for the complexity analyses below we will use that a C$^p$NOT gate requires $\mathcal{O} \left (p^2 \right )$ CNOT gates and one ancilla qubit. 

We would like to stress here that there might be cheaper methods for implementing the comparison operation and/or a multi-controlled NOT operation possible as this is a topic of active research. This will only decrease the total cost of our method and it will not change the conclusion presented in Section \ref{ssec:comp_soa_q} of our method having the lowest complexity of any known CQBM. 

\subsection{Complexity of our CQBM solver}\label{ssec:multi not impl}
In this section we give a detailed analysis of the costs of our CQBM solver by first analyzing the different steps separately and, subsequently, combining these results to find the total complexity of the algorithm. 

\subsubsection{Complexity of quantum state preparation}\label{sssec:stateprep}
In order to run the CQBM solver, we first need to prepare the input state to represent the desired particle distribution and velocities, at the initial time $t=0$.
As it is known that the preparation of arbitrary states can become exponentially expensive \cite{Plesch2011, mikeandike}, we restrict ourselves for the numerical results presented in these notes to equidistributed particles and mention a few publications that describe efficient preparation procedures for particular distributions.
Preparing general sparse quantum states is known to be possible with circuit depth $\Theta \log \left (ns\right )$ using $\mathcal{O}\left ( ns \log(s) \right)$, where $s $ the number of nonzero entries of the quantum state \cite{Zhang2022}. Since our input quantum states is typically highly sparse, this is a good result for the efficient state preparation for our quantum states.\\
Another promising result for the input state preparation is given by Grover and Rudolph \cite{Grover2002}, who present an efficient process for preparing quantum states that form a discrete approximation of an efficiently integrable probability density function, such as log-concave distribution functions. Since the Maxwell-Boltzmann distribution as described in Section \ref{sec:collisionless boltzmann equation} is log-concave our future research will focus on exploiting this property for the efficient quantum state preparation.

\subsubsection{Complexity of the convection step}\label{sssec:compl conv step}
The quantum convection step consists of a streaming operation applied to the particles traveling at a pre-determined velocity. In order to implement this we first apply a $\text{C}^p\text{NOT}$ gate between the velocity qubits in each dimension $i$ and the $a_{u,i}$ ancilla qubits. Each $\text{C}^p\text{NOT}$ gate has a total of $p=n_{v_i}-1$ control qubits. We implement this operation for the particles whose velocity magnitudes are marked for advance in that particular timestep. The amount of velocity magnitudes being advanced in each timestep $n^v_t$ can trivially be capped by $N_v$, but in practice will be equal to $1$ in $89.6\%$ of the cases. We found this percentage by performing numerical analysis for $N_v$ up to 1024. 
Subsequently, controlled on the $a_{u,i}$ qubits the increase operator is applied.
Our increase operator consists of a QFT operation applied to the $n_{g_i}$ qubits in each dimension $i$, followed by a phase shift on the $n_{g_i}$ qubits and a final QFT$^\dagger$ operation. In each dimension $i$, a QFT operation requires $\mathcal{O}(n_{g_i}^2)$ two-qubit operations \cite{Musk2020}, the phase shifts applied to the $n_{g_i}$ qubits is controlled on the $a_{u,i}$ ancilla qubits and so this operation costs $\mathcal{O}(n_{g_i}^2)$ two-qubit operations. 
Therefore our increase operator only requires $\mathcal{O}(n_{g_i}^2)$ CNOT operations, controlled on the $a_{u,i}$ qubits in each dimension $i$.
This allows us to cap the complexity of our increase operator by $\mathcal{O}(d n_{g_\text{max}}^2)$ CNOT operations.
%which is an exponential reduction in case there are no ancilla qubits available, otherwise our created circuits still require far less CNOT gates and requires $\mathcal{O}(n)$ less ancilla qubits; See Figure \ref{fig:com_with_ancillae} for a comparison of the amount of CNOT gates required and circuit depth of our proposed incrementor with the original one. . 
% For QFT there is a rather straightforward method to use a afschatting instead of the full QFT. When playing around with this we get the best result when ... 

In total that leaves us with a complexity of the convection step of $\mathcal{O} \left ( d n^v_t \right )$ in the amount of C$^p$NOT operations with $p=n_{v_\text{max}}-1$ plus $\mathcal{O} \left ( d n_{g_\text{max}}^2 \right )$ CNOT operations.\\

Using the complexity analysis of the multi-controlled NOT operations provided in Subsection \ref{ssec:multinot compl}, we get an overall complexity of $$\mathcal{O} \left ( d n^v_t n_{v_\text{max}}^2 + d n_{g_\text{max}}^2 \right )$$ in the amount of CNOT gates for the convection step.
%Here, clearly the dominant factor are the $\mathcal{O} \left ( d n^v_t \right )$ multi-controlled NOT gates. This is quite efficient since in practice $d\leq 3$ will be used and $n_v^t=1$ in $\approx90\%$ of the timesteps. 

\subsubsection{Complexity of the specular reflection step}
The specular reflection step for each wall starts with two quantum comparison operations, applied to the qubits encoding the location in a single dimension with the $a_{l}$ and $a_{u}$ qubits storing the result. This leads to a complexity of $\mathcal{O} \left ( n_{g_\text{max}}^2 \right) $ CNOT gates when using the Draper adder \cite{Draper1998} to implement the comparison operations. 

Subsequently a multi-controlled NOT gate applied to the ancilla qubits $a_{o,i}$, $i \in \{1,\dots,d\}$, controlled on the position in the dimension the considered wall reflects, the $2(d-1)$ $a_{u}$, $a_{l}$ ancillae and one $a_{u,i}$ qubits and one directional qubit $v_{\text{dir},i}$ is applied. Therefore the second step of the specular reflection for each wall consists of a $\text{C}^{p}\text{NOT}$ operations with $p \leq n_{g_\text{max}} +2\left ( d-1 \right ) +2 =  n_{g_\text{max}} +2d$. 

Since the two steps above are applied to each wall we get a total complexity of $\mathcal{O} \left ( n_w n_{g_\text{max}}^2 \right) $ CNOT gates and $\mathcal{O}\left ( n_w \right )$ $\text{C}^{p}\text{NOT}$ operations with $p \leq n_{g_\text{max}} +2d$ and $n_w$ the total number of walls. 

Then, controlled on the $a_{o,i}$ qubits we perform a NOT operation on the $v_{\text{dir},i}$ qubits. This operation consists of $d$ CNOT gates, and so the complexity is $\mathcal{O}\left ( d\right )$ CNOT gates. This is followed by a single incrementation operation applied to the positional qubits controlled on the $a_{o,i}$ in each dimension. 
The incrementation operation again has complexity $\mathcal{O} \left (d n_{g_\text{max}}^2 \right )$ CNOT operations. 

Subsequently, we need to reset the $a_{o,i}$ ancillae. 
This is done for each wall by again first performing two quantum comparison operations, followed by a multi-controlled NOT gate targeting the $a_{o,i}$ ancillae controlled on the qubits encoding the position in a single dimension in combination with the $2(d-1)$ $a_{u}$, $a_{l}$ ancillae, one $a_{u,i}$ qubit and one directional qubit $v_{\text{dir},i}$. This leaves us with a total complexity of again $\mathcal{O} \left ( n_w n_{g_\text{max}}^2 \right) $ CNOT gates and $\mathcal{O}\left ( n_w \right )$ $\text{C}^{p}\text{NOT}$ operations with $p \leq n_{g_\text{max}} +2d$.

Only for the resetting of the $a_{o,i}$ qubits we need to take into account the special rules for resetting the corner cases by applying a multi-controlled NOT gate targeting an ancilla $a_{o,i}$, controlled on the position as well as the $a_{u,i}$ qubits and the direction qubits $v_{\text{dir},i}$. Therefore the complexity of resetting an ancilla $a_{o,i}$ for a corner case becomes $\mathcal{O}\left (  1 \right )$ $\text{C}^{p}\text{NOT}$ operations with $p \leq n_g +2d$. Since the amount of corner cases is linear in the amount of walls the total complexity becomes $\mathcal{O}\left (  n_w  \right )$ $\text{C}^{p}\text{NOT}$
operations with $p \leq n_g +2d$.

At the end of the specular reflection step, we reset the $a_{u,i}$ ancilla qubits in each dimension $i$, using a $\text{C}^p\text{NOT}$ gate with $p=n_{v_i}-1$. 
As explained in Section \ref{sssec:compl conv step} the total costs of this operation amount to $\mathcal{O} \left ( d n^v_t \right )$ C$^p$NOT operations with $p=n_{v_\text{max}}-1$. \\

In total, the complexity of the reflection step is equal to
$\mathcal{O} \left ( n_w n_{g_\text{max}}^2 \right) $ CNOT gates plus $\mathcal{O}\left ( n_w \right )$ $\text{C}^{p}\text{NOT}$ operations with $p \leq n_{g_\text{max}} +2d $ plus $\mathcal{O}\left (  n_w  \right )$ $\text{C}^{p}\text{NOT}$
operations with $p \leq n_g +2d$ plus $\mathcal{O}\left ( d\right )$ CNOT gates plus $\mathcal{O} \left (d n_{g_\text{max}}^2 \right )$ CNOT gates. 

When combining these we get a total complexity of $\mathcal{O} \left ( (n_w + d)  n_{g_\text{max}}^2  \right)  $ CNOT gates
plus $\mathcal{O}\left (  n_w  \right )$ $\text{C}^{p}\text{NOT}$
operations with $p \leq n_g +2d$, for the complexity of the reflection step. 

Translating these results in terms of multi-controlled NOT complexities into single-controlled NOT complexities using the results of Subsection \ref{ssec:multinot compl}, we get the following total complexity 
$$\mathcal{O}\left (  n_w \left ( n_g +2d \right)^2 + (n_w + d)  n_{g_\text{max}}^2 \right  ) = \mathcal{O}\left (  n_w n_g^2+ (n_w + d)  n_{g_\text{max}}^2 \right ) = \mathcal{O}\left (  n_w n_g^2 \right )$$ CNOT gates.

\subsubsection{Total CQBM complexity}
Combining the complexities of the convection and specular reflection step as detailed above, we obtain that the total complexity of the algorithm per timestep amounts to $\mathcal{O} \left ( (n_w + d)  n_{g_\text{max}}^2  \right)  $ CNOT gates and $\mathcal{O} \left ( n_w \right )$ $\text{C}^{p}\text{NOT}$ gates with $p \leq n_g +2d$ and $\mathcal{O} \left ( d n^v_t \right )$ in the amount of C$^p$NOT operations with $p=n_{v_\text{max}}-1$.  \\

Using the multi-controlled NOT decomposition as before, this leads to a total complexity of 
$$\mathcal{O}\left (  n_w n_g^2 + d n^v_t n_{v_\text{max}}^2 + d n_{g_\text{max}}^2 \right ) =\mathcal{O}\left (  n_w n_g^2 + d n^v_t n_{v_\text{max}}^2 \right )$$ CNOT gates per time step.


\subsection{Complexity of alternative CQBM implementations}\label{ssec:comp_soa_q}
An alternative collisionless Quantum Boltzmann solver presented in the literature is the one by Todorova and Steijl \cite{Todorova2020}. Though not being fully fail-safe in the specular reflection step, the authors present some complexity analysis which led us to perform a rigorous comparison of the complexities of both approaches.

As reported in their paper, Todorova and Steijl conclude a complexity in terms of $\text{C}^p\text{NOT}$ gates of $\mathcal{O}\left ( d N_v \log_2 \left ( D / h  \right ) \right )$ for the streaming step per timestep with $p=n_{g_\text{max}} + n_{v_\text{max}}$. The complexity of their specular-reflection boundary condition, again, quantified in terms of C$^p$NOT gates is $\mathcal{O}\left ( d N_v \log_2 \left ( D / h  \right ) \right )$. Here the authors do not provide an explicit estimation for $p$, but since they control on positions in all $d$ spatial dimensions and streaming speeds combined, we derive a complexity of $p = n_{g} +n_{v_\text{max}}$.

Since assumptions and parameters in the paper by Todorova and Steijl differ from ours, it is not immediately clear how to compare their complexity analysis with ours on a fair basis. 
Rewriting their complexity result in terms of the parameters adopted in our analysis yields $\mathcal{O}\left ( d N_v \log_2 \left ( D / h  \right ) \right ) = \mathcal{O}\left  (d N_v n_{g_\text{max}} \right ) $ C$^p$NOT operations with $p \leq n_{g} +n_{v_\text{max}}$.
Finally, rewriting their complexity result in terms of single-controlled CNOT rather than C$^p$NOT gates, yields a total complexity of $\mathcal{O}\left ( d N_v n_{g_\text{max}} (n_{g} +n_{v_\text{max}})^2 \right )$.

Another difference that makes our complexity hard to compare with the complexity of the algorithm of Todorova and Steijl is that the aforementioned authors do not adopt the variable $n_v^t$, and instead use the maximum value $N_v$. For ease of comparison we will replace $n_v^t$ with $N_v$ in our estimates which is possible as $n_v^t \leq N_v$ holds trivially.
Finally, Todorova and Steijl also do not take the amount of walls into account explicitly, but instead seem to consider it a constant, we need to relax this parameter in our complexity estimation as well. With all the aforementioned changes in place we can rewrite our leading complexity terms as $\mathcal{O}\left ( n_g^2 + d N_v n_{v_\text{max}}^2 \right )$ CNOT operations, which is significantly cheaper than the CQBM approach by Todorova and Steijl having a total CNOT complexity of $\mathcal{O}\left ( d N_v n_{g_\text{max}} (n_{g} +n_{v_\text{max}})^2 \right )$.

\subsection{Complexity comparison of incrementation operations}\label{ssec:complexity incrementation}
One of the improvements we proposed in the material of these notes with respect to other known methods \cite{Todorova2020, Budinski2020, Budinski2021} is the use of the Quantum Draper Adder \cite{Draper1998} inspired approach to implement the quantum incrementation operation, leading to a cheaper quantum primitive for the streaming operation. In this section we provide a comparison of the complexity of the incrementation method with our QDA inspired approach compared to the approaches implemented in \cite{Todorova2020, Budinski2020, Budinski2021}. 

Consider the QDA inspired incrementation operation applied to the qubits $g_i$ encoding the position in dimension $i$, this consists of a QFT followed by a phase shift gate and finally a QFT$^\dagger$ operation. 
The cost of implementing the QFT (QFT$^\dagger$) operation is $n_{g_i}^2 + \frac{1}{2}n_{g_i}$ CNOT operations as shown in Section \ref{ssec:multinot compl}. Therefore our incrementation operation can be applied at a total cost of $2 n_{g_i}^2 + n_{g_i}$ CNOT operations. 
%This leads to a total complexity of $$\mathcal{O}\left ( n_{g_i}^2 \right )$$ CNOT operations associated with the QDA inspired incrementation approach. \\
%Note that the constant here is relatively small and for $n_{g_i} \geq 4$ we use less than or equal to $2 n_{g_i}^2$ CNOT gates\footnote{A QFT operation on $n_{g_i}$ qubit requires $\frac{n_{g_i}(n_{g_i}-1)}{2}$ controlled phase shift operations followed by $\lfloor \frac{n_{g_i}}{2} \rfloor$ swap operations. A swap operation can be implemented using 3 CNOT operations \cite{mikeandike} and a controlled phase shift operation can be implemented using a single CNOT operation \cite{Qiskit}.}.  

The cost of the quantum incrementation operations proposed in \cite{Todorova2020, Budinski2020, Budinski2021}\footnote{Note that in the papers \cite{Budinski2020, Budinski2021} the primitive for the incrementation operation is not explicitly given, we have deduced the implementation of the incrementation in our method using the Figures provided.} consists of $\sum_{p=0}^{n_{g_i}-1} \text{C}^p\text{NOT}  $ gates. Only when assuming that the costs of implementing a C$^p$NOT gate consists of $\mathcal{O}\left ( p\right )$ NOT gates, we find that the total costs of the incrementation operation implemented in these papers have the same order of magnitude with respect to the amount of CNOT gates required as our QDA inspired incrementation method. To the best of our knowledge, however, a decomposition of a C$^p$NOT gate that requires $\mathcal{O}\left ( p\right )$ NOT gates, either requires an extra $p-1$ ancilla qubits or requires a single ancilla qubit and $c$ CNOT operations with a large constant $c$, making our method more efficient \ref{ssec:multinot compl}. 
